\section{Adding classical measurement error}
\label{sec:measurement-error}

The previous section concerns a phenotype measured without error. We now
separate that phenotype from the recorded outcome. This is a specific
measurement model: adding a residual and calling it noise is not enough
to justify multiplying every kinship correlation by the same factor.

\subsection{The observed outcome and the meaning of the attenuation factor}

Retain the standardized genetic and environmental components $A_i,E_i$,
but let the recorded outcome be
\begin{equation}
 Y_i=aA_i+eE_i+uU_i,\qquad
 X_i^+=\begin{pmatrix}A_i\\E_i\\U_i\end{pmatrix},\qquad
 w=\begin{pmatrix}a\\e\\u\end{pmatrix},\qquad Y_i=w^{\mathsf T}X_i^+.
 \label{eq:me-outcome}
\end{equation}
$U_i$ represents pure measurement error, with mean zero and variance
one. Assume $a,e>0$, $u\geq0$, and, for every pair of people $i,j$,
\begin{equation}
 \operatorname{Cov}(U_i,A_j)=\operatorname{Cov}(U_i,E_j)=0,\qquad
 \operatorname{Cov}(U_i,U_j)=0\quad(i\ne j).
 \label{eq:me-orthogonality}
\end{equation}
Thus error is uncorrelated with both true components within the person
and across people, and with other people's errors. Its variance within
the person is one, not zero. The original within-person condition
$\operatorname{Cov}(A_i,E_i)=0$ is retained; the nonzero cross-person
genetic--environmental covariances derived earlier are also retained.
These orthogonality conditions suffice for the moment calculations.
Independent errors, independent of all true family components, are a
stronger sufficient specification used in the simulations.

Standardize each recorded outcome to unit variance. The three
within-person components are uncorrelated, so
\[
 \operatorname{Var}(Y_i)=a^2+e^2+u^2=1.
\]
Define the attenuation factor and noise share as
\begin{equation}
 \eta:=u^2,\qquad
 \theta:=1-\eta=a^2+e^2>0.
 \label{eq:me-shares}
\end{equation}
Here $\eta$ is the fraction of observed variance attributable to
measurement error, whereas $\theta$ is the fraction attributable to
the true phenotype. This distinction matters for translating Clark's
notation: his $\theta$ is the \emph{attenuation factor}
\citep[p.~42 and Table~A1]{Clark2026}, so
\[
 \boxed{\theta_{\rm Clark}=\theta=1-\eta.}
\]
Thus $1-\theta=\eta$ is the noise share.

The standardized latent phenotype is
\begin{equation}
 P_i=\frac{aA_i+eE_i}{\sqrt{\theta}}
     =hA_i+e_P E_i,\qquad
 h=\frac{a}{\sqrt{\theta}},\quad
 e_P=\frac{e}{\sqrt{\theta}}=\sqrt{1-h^2}.
 \label{eq:me-latent}
\end{equation}
Consequently
\[
 Y_i=\sqrt{\theta}P_i+\sqrt{\eta}U_i,\qquad
 h^2=\frac{a^2}{\theta},\quad
 a^2=\theta h^2,\quad e^2=\theta(1-h^2).
\]
$h^2$ is heritability of the latent phenotype; $a^2$ is the additive
genetic share of variance of the recorded outcome. They coincide only
when $\eta=0$. The symbol $e_P$ denotes the environmental loading in
the error-free model; $e$ denotes its loading on the observed scale.

Assume PH and own-person conditional linearity for this same latent
$P_i$, together with the earlier transmission and pedigree assumptions.
Spouses match on $P$, not on the recording error $U$. Therefore
\begin{equation}
 \rho_A=h^2\rho=\frac{a^2}{\theta}\rho,\qquad
 \rho:=\operatorname{Corr}(P_M,P_F).
 \label{eq:me-spouse-restriction}
\end{equation}
No new transmission term is introduced: $A_j=(A_M+A_F)/2+S_j$ and
$\operatorname{Var}(S_j)=(1-h^2\rho)/2$ as before. $U$ changes what is
recorded, not what is inherited or how mates are selected.

\subsection{Why the same attenuation factor appears for individual relatives}
\label{sec:me-attenuation}

For distinct people, expanding the two observed outcomes gives
\begin{align}
 \operatorname{Cov}(Y_i,Y_j)
 &=\theta\operatorname{Cov}(P_i,P_j)
   +\sqrt{\theta\eta}\operatorname{Cov}(P_i,U_j)\notag\\
 &\quad+\sqrt{\theta\eta}\operatorname{Cov}(U_i,P_j)
   +\eta\operatorname{Cov}(U_i,U_j).
 \label{eq:me-expansion}
\end{align}
For example,
$\operatorname{Cov}(P_i,U_j)
=h\operatorname{Cov}(A_i,U_j)+e_P\operatorname{Cov}(E_i,U_j)=0$
by \eqref{eq:me-orthogonality}; the other cross-term vanishes in the
same way, and the error--error covariance is zero. Since both observed
and latent individual phenotypes have unit variance,
\begin{equation}
 \boxed{\operatorname{Corr}(Y_i,Y_j)
   =\theta\operatorname{Corr}(P_i,P_j)
   =(1-\eta)\operatorname{Corr}(P_i,P_j),\quad i\ne j.}
 \label{eq:me-attenuation}
\end{equation}
This is the source of the common multiplier in the individual-relative
rows of Clark's Table~A1. On the unstandardized signal-plus-error scale,
measurement error adds variance without adding covariance; the
correlation is therefore reduced. Attenuation occurs at the two observed
endpoints, not anew at each intervening generation.

In particular, the observed spouse correlation is
$s_Y:=\operatorname{Corr}(Y_M,Y_F)=\theta\rho$. Thus the latent
restriction may also be written
\[
 \rho_A=\frac{a^2s_Y}{\theta^2};
\]
it is not generally $a^2s_Y$. Substituting observed quantities into the
error-free restriction without the scaling factors is incorrect.
The complete spouse block for the standardized components is
\[
 \operatorname{Cov}(X_M^+,X_F^+)
 =\rho\begin{pmatrix}
 h^2&he_P&0\\he_P&e_P^2&0\\0&0&0
 \end{pmatrix}.
\]
Multiplying on the left by $w^{\mathsf T}$ and on the right by $w$
gives $\rho(ah+ee_P)^2
=\rho\{\sqrt{\theta}(h^2+e_P^2)\}^2=\theta\rho$, confirming the result.

\subsection{The adjusted moment table}
\label{sec:me-table}

Apply \eqref{eq:me-attenuation} to each moment already derived. For
example,
\[
 r_{PO}^{Y}=\theta\frac{h^2(1+\rho)}{2},\qquad
 r_S^{Y}=\theta h^2k,\qquad
 k:=\frac{1+h^2\rho}{2}.
\]
Here $r^Y$ denotes a correlation of recorded outcomes. The analogous
substitution gives the grandparent, avuncular, and first-cousin rows
in Table~\ref{tab:measurement-moments}.

The more distant collateral rows follow from the same pedigree
extension, rather than from adding further error factors.
For two collateral relatives $i,j$ whose cross-component environmental
covariances vanish, equation~\eqref{eq:model-child-mean} implies
\[
 \operatorname{Cov}(A_{c(i)},A_j)
 =k\operatorname{Cov}(A_i,A_j)+d\operatorname{Cov}(E_i,A_j)
 =k\operatorname{Cov}(A_i,A_j).
\]
Fresh offspring environments preserve the zero environmental
cross-covariances. Descending one side therefore multiplies genetic
covariance by $k$; descending both sides multiplies it by $k^2$,
using conditionally independent outside mates and births as in the
first-cousin derivation. Starting with sibling covariance $k$ gives
$k^3,k^5,k^7,k^9$ for first through fourth cousins. Removing a cousin
relationship by $t$ generations on one side supplies another $k^t$.
The observed phenotype correlation remains $\theta h^2$ times that
genetic covariance.

\begin{table}[htbp]
\centering
\caption{Correlations with classical measurement error:
$\theta=1-\eta$, $k=(1+h^2\rho)/2$}
\label{tab:measurement-moments}
\begin{tabular}{lccc}
\hline
Relationship & $n$ & $D$ & Observed correlation $r^Y$\\
\hline
Spouses & --- & --- & $\theta\rho$\\
Parent--child & 1 & 1 & $\theta h^2(1+\rho)/2$\\
Full siblings & 1 & 0 & $\theta h^2k$\\
Grandparent--grandchild & 2 & 1 & $\theta h^2(1+\rho)k/2$\\
Aunt/uncle--niece/nephew & 2 & 0 & $\theta h^2k^2$\\
First cousins & 3 & 0 & $\theta h^2k^3$\\
Second cousins & 5 & 0 & $\theta h^2k^5$\\
Third cousins & 7 & 0 & $\theta h^2k^7$\\
Fourth cousins & 9 & 0 & $\theta h^2k^9$\\
Degree-$c$ cousins, $t$ removed & $2c+1+t$ & 0
 & $\theta h^2k^{2c+1+t}$\\
\hline
\end{tabular}
\end{table}
In the final row $c\geq1$ and $t\geq0$. The columns $n,D$ specify the
regression representation below; the spouse row is a separate moment.
The single-parent and single-grandparent rows, and the collateral rows,
agree with Clark's Table~A1 after using $\theta$ as the common attenuation factor,
$r=\rho$, and $m=h^2\rho$. Averages of relatives require a different
normalization and are treated separately below.

\subsection{What the appendix regression can identify}
\label{sec:me-identification}

For the nonspousal rows, both the collateral and lineal formulas can
be written
\begin{equation}
 r^Y_{n,D}
 =\theta h^2 k^n
  \left(\frac{1+\rho}{1+h^2\rho}\right)^D,\qquad D\in\{0,1\}.
 \label{eq:me-regression-moment}
\end{equation}
For example, $n=1,D=1$ cancels $1+h^2\rho$ from $k$, leaving the
parent--child formula. $n=2,D=1$ leaves the extra factor $k$ for a
grandparent. For positive population correlations, taking logarithms gives
\begin{align}
 \log r^Y_{n,D}&=\alpha+\beta n+\gamma D,\notag\\
 \alpha&=\log(\theta h^2),\quad
 \beta=\log\left(\frac{1+\rho_A}{2}\right),\quad
 \gamma=\log\left(\frac{1+\rho}{1+\rho_A}\right).
 \label{eq:me-log-regression}
\end{align}
This is the structure of Clark's equations~A1--A3; Greek letters here
avoid confusing the regression intercept with the loading $a$.

With sufficient variation to identify all three regression coefficients,
invert them in sequence:
\begin{align}
 k&=\exp(\beta),&
 \rho_A&=2\exp(\beta)-1,\notag\\
 1+\rho&=(1+\rho_A)\exp(\gamma),&
 \rho&=2\exp(\beta+\gamma)-1,\notag\\
 h^2&=\frac{\rho_A}{\rho},&
 \theta&=\frac{\exp(\alpha)}{h^2},&
 \boxed{\eta=1-\frac{\exp(\alpha)}{h^2}.}
 \label{eq:me-inversion}
\end{align}
The crucial third line invokes the \emph{derived PH restriction}
$\rho_A=h^2\rho$. Without it, the intercept determines only
$\theta h^2=a^2$, not the separate noise and heritability shares.
Thus estimating noise by this route is possible in principle, but
its interpretation depends on the same substantive model being tested.
The exponentiated intercept is the observed genetic variance share $a^2$;
calling it latent heritability $h^2$ would be incorrect when noise is present.

The separation works because genuine $E$ enters the phenotype on which
spouses match, whereas $U$ does not. Their effects on the pattern of
parent--child and sibling resemblance therefore differ under the model.
Without that assumed distinction, two independent, untransmitted sources
of variation need not be distinguishable from these data. The recovered
noise share is consequently not an independent validation of the mating
assumptions used to identify it.

Two limits are important. With \emph{collateral relatives alone}, $D=0$:
the regression identifies $\theta h^2$ and $\rho_A$, but not $\rho$,
$h^2$, and $\eta$ separately. Even after imposing PH, different pairs
$(\theta,h^2)$ can have the same product, with
$\rho=\rho_A/h^2$, provided all parameters remain admissible.
Also, under random mating $\rho=\rho_A=0$, $\gamma=0$ for all $h^2$;
the ratio $\rho_A/\rho$ cannot identify heritability. The preceding
inversion therefore presumes positive assortment and an identified
lineal--collateral contrast. Identification can be weak near random
mating.

Estimation must enforce $0<h^2<1$, $0<\rho<1$, and
$0<\theta\leq1$. Arbitrary fitted log-regression coefficients need
not imply admissible shares. One can instead fit the three structural
parameters $(h^2,\rho,\eta)$ jointly to the moment table, imposing
the constraints directly. The spouse moment $\theta\rho$, if available,
provides an additional check on the predictions recovered from the
nonspousal regression.

Equation~\eqref{eq:me-log-regression} is exact for population moments;
using logarithms of estimated correlations introduces sampling issues.
Logs require positive estimates, and
$\operatorname{SE}(\log\widehat r)\simeq
\operatorname{SE}(\widehat r)/r$ by the delta method. For inverse-variance weighting of the log regression, the relevant
uncertainty is on the log scale. Inference must also account for
dependence between estimates from overlapping families. A small
residual sum of squares does not by itself verify the measurement or
transmission assumptions that identify $\eta$.

\subsection{Averaging parents changes the correlation}
\label{sec:me-averages}

The common multiplier $\theta$ above applies to pairs of \emph{individuals}.
It cannot be carried over mechanically to the correlation of a child
with an average of relatives. For the observed midparent value
$\overline Y=(Y_M+Y_F)/2$,
\begin{align*}
 \operatorname{Var}(\overline Y)
 &=\tfrac14\{1+1+2\operatorname{Cov}(Y_M,Y_F)\}
   =\tfrac12(1+\theta\rho),\\
 \operatorname{Cov}(Y_j,\overline Y)
 &=\tfrac12\{\operatorname{Cov}(Y_j,Y_M)
             +\operatorname{Cov}(Y_j,Y_F)\}\\
 &=\theta h^2(1+\rho)/2.
\end{align*}
Since $\operatorname{Var}(Y_j)=1$, the correlation is
\begin{equation}
 \boxed{\operatorname{Corr}(Y_j,\overline Y)
 =\frac{\theta h^2(1+\rho)}
        {\sqrt{2(1+\theta\rho)}}.}
 \label{eq:me-midparent}
\end{equation}
Even without measurement error this is
$h^2\sqrt{(1+\rho)/2}$, not $h^2$. The familiar $h^2$ is instead the
error-free \emph{regression slope} on the midparent phenotype.
With error, that slope becomes
\[
 \frac{\operatorname{Cov}(Y_j,\overline Y)}
      {\operatorname{Var}(\overline Y)}
 =\frac{\theta h^2(1+\rho)}{1+\theta\rho},
\]
which is generally not $\theta h^2$ either.

Clark's Table~A1 labels the average-parent entry a correlation and gives
$\theta h^2$. Interpreted literally as the correlation with the arithmetic
mean of the parents' outcomes, that entry does not follow from this model.
It needs a different definition or a correction. The average-grandparent
entry requires the same care: its denominator depends on the joint
covariance matrix of the four grandparents. Averaging individual-relative
correlations is not the same as correlating with an average.

\subsection{Substantive content of the common-noise correction}
\label{sec:me-meaning}

The adjustment assumes more than imperfect measurement. First, each
person's error must be uncorrelated with \emph{every} relevant relative's
true components and error. A shared informant, shared record linkage
error, or common error in assigning status to an occupation can violate
that requirement. For example, if the two error--signal terms in
\eqref{eq:me-expansion} remain zero but
$\operatorname{Cov}(U_i,U_j)=c_{ij}$, then
\[
 r^Y_{ij}=\theta r^P_{ij}+\eta c_{ij}.
\]
This is not a uniform attenuation of the original formula.

Second, the common multiplier presumes the same signal share for the
people entering each correlation. If instead
$Y_i=\sqrt{\theta_i}P_i+\sqrt{1-\theta_i}U_i$ with the analogous
orthogonality conditions, expansion gives
\[
 r^Y_{ij}=\sqrt{\theta_i\theta_j}\,r^P_{ij}.
\]
Different reliability by sex, age, cohort, outcome definition, or kinship
sample can therefore prevent cancellation in ratios. Those restrictions
must hold in the actual selected samples, not merely in a hypothetical
unselected population.

Third, the latent phenotype must be the same trait whose components
drive transmission and assortment. If an observed characteristic
affects partner choice, the part of that characteristic labeled error
relative to an unobserved trait cannot automatically be treated as
irrelevant noise. Nor does an imperfect proxy for intelligence
automatically satisfy classical measurement-error assumptions for
intelligence. A substantive difference between traits may contain
systematic genetic or environmental variation.

Finally, pure measurement error is distinct from the true environmental
component $E$. Independent, untransmitted environmental variation can
also reduce familial resemblance, and its separation from recording
error here depends on the modeled assortment and transmission
restrictions. Correcting attenuation does not remove the exclusion
of shared upbringing, validate PH, or establish a genetic explanation
for a close fit. Under a common positive $\theta$ the parent--child
minus sibling gap retains its sign:
\[
 r^Y_{PO}-r^Y_S
 =\theta h^2\rho(1-h^2)/2>0.
\]
Allowing arbitrary relationship-specific errors would change that test,
but it would also change the model being fitted.

\paragraph{Verification.}
The measurement-error verification script checks the covariance
expansion, scaling, regression inversion, and midparent normalization
symbolically. It then adds fresh measurement errors to generated
pedigrees and checks the predicted individual-relative and midparent
moments numerically. All 71 new symbolic checks passed, along with the 46 baseline checks. Across 13 parameter settings, 2.6 million simulated pedigrees passed 156 individual-relative and 26 midparent checks, including distant cousins and relatives once removed. It also checks counterexamples with unequal
reliability and correlated measurement errors. These calculations
verify the stated model; they do not verify its empirical assumptions.
